\chapter*{Answer Key}
\addcontentsline{toc}{chapter}{Answer Key}

\section*{Chapter 1: Introduction}

\begin{enumerate}[leftmargin=*]
\item Preserve life, prevent the condition from worsening, and promote recovery.
\item Immediate recognition of cardiac arrest and activation of the emergency response system; early CPR with emphasis on chest compressions; rapid defibrillation when available; advanced life support by healthcare professionals; and post-cardiac arrest care.
\item The helper acts voluntarily (not expecting payment); provides care within their level of training; does not act recklessly or with gross negligence; and obtains consent when possible (implied consent is assumed for unconscious persons).
\item Check scene safety. Your safety comes first -- approaching an unsafe scene makes you another casualty and reduces the help available.
\item Call emergency services (or direct a bystander to) and begin CPR immediately, starting with chest compressions. Recognition and activation come first, then compressions.
\end{enumerate}

\section*{Chapter 2: Basic First Aid Principles}

\begin{enumerate}[leftmargin=*]
\item Wear disposable gloves; avoid direct skin-to-skin contact with body fluids; wash hands before and after care; dispose of contaminated materials and sharps safely.
\item With the head-tilt/chin-lift opening the airway, look, listen, and feel for normal breathing for 5--10 seconds. An unresponsive person who is not breathing normally needs CPR.
\item Use it when a person is unresponsive but breathing normally, to keep the airway open and allow fluids to drain. Avoid it when spinal injury is suspected (unless the person is vomiting or not breathing), and modify it for pregnant persons (left side) and infants.
\item At least 5 minutes of continuous direct pressure before checking.
\item They damage tissue and delay healing. Clean the wound gently with sterile saline or clean water instead.
\end{enumerate}

\section*{Chapter 3: Assessment}

\begin{enumerate}[leftmargin=*]
\item Scene safety, responsiveness, call for help, then Airway, Breathing, and Circulation (ABC). Interrupt the survey immediately to treat any life-threatening problem you find.
\item Place them in the recovery position, monitor breathing continuously, and keep emergency services on the way. Stay until help arrives.
\item Signs and symptoms; Allergies; Medications; Past medical history; Last oral intake; Events leading up to the incident. It provides a structured history for responders.
\item AVPU (Alert, Verbal response, Pain response, Unresponsive) is a quick way to record level of consciousness repeatedly, so any deterioration is detected early.
\item Signs include pale, cool, clammy skin; a rapid weak pulse; rapid shallow breathing; weakness or altered mental status. First steps: call emergency services, lay the person flat, keep them warm, and manage the cause. (Elevate the legs only for non-traumatic causes, and never if it causes pain.)
\end{enumerate}

\section*{Chapter 4: Bleeding Control}

\begin{enumerate}[leftmargin=*]
\item Arterial bleeding is bright red and spurts with the pulse; venous bleeding is darker and flows steadily. Distinguishing them helps judge severity -- arterial bleeding is more likely to be life-threatening.
\item Removing a soaked dressing tears away the forming clot and restarts serious bleeding. Add additional dressings on top instead.
\item A tourniquet is appropriate for life-threatening limb bleeding not controlled by direct pressure. Always note the time it was applied, place it 2--3 inches above the wound (not over a joint), tighten until bleeding stops, and never loosen it.
\item Do NOT remove it. The object may be plugging the wound and controlling bleeding. Stabilize it in place with bulky dressings around it, control bleeding around it, and seek emergency medical care.
\item Sit the person leaning slightly forward and pinch the soft part of the nose with firm pressure for at least 10 minutes continuously, breathing through the mouth. Do not tilt the head back.
\end{enumerate}

\section*{Chapter 5: CPR and Airway Management}

\begin{enumerate}[leftmargin=*]
\item Place the heel of one hand on the center of the chest, lower half of the sternum, other hand on top, fingers interlocked. Compress 2--2.4 inches (5--6 cm) at a rate of 100--120 per minute, allowing full chest recoil.
\item 30 compressions to 2 breaths for a single rescuer on an adult. The same 30:2 ratio applies for a single rescuer on an infant (with compression depth about 1.5 inches/4 cm).
\item Head-tilt/chin-lift can displace an injured cervical spine and cause paralysis. Use the jaw-thrust maneuver instead, maintaining manual in-line stabilization.
\item Move a child or infant during CPR only if the scene is unsafe (fire, falling debris, traffic). Move them as a whole unit supporting the head and neck -- a young child or infant can be carried in the forearms (head cradled) while CPR continues. If the scene is safe, leave them in place.
\item A shockable rhythm (ventricular fibrillation or pulseless ventricular tachycardia) means the AED will advise a shock; deliver it, then resume CPR immediately for 2 minutes before re-analyzing. If no shock is advised, resume CPR immediately and let the AED re-analyze every 2 minutes.
\item For an infant: give 5 back blows between the shoulder blades, then 5 chest thrusts, alternating until the object is dislodged (never abdominal thrusts). For an adult: give abdominal thrusts (Heimlich maneuver) until the object is expelled or the person becomes unresponsive.
\end{enumerate}

\section*{Chapter 6: Medical Emergencies}

\begin{enumerate}[leftmargin=*]
\item Shortness of breath, nausea or vomiting, fatigue or weakness, pain in the jaw, back, or between the shoulder blades, and cold sweats -- sometimes with little or no chest pressure.
\item Time-based treatments such as clot-busting drugs must be given within a narrow window from the moment symptoms began. Recording onset time determines eligibility and guides treatment.
\item Call emergency services immediately (a seizure lasting longer than 5 minutes is a medical emergency). Protect the person from injury, time the seizure, do not restrain them or put anything in their mouth, and place them in the recovery position once it stops.
\item When the person has an aspirin allergy, has been advised by their doctor against it, has signs of a bleeding condition such as stroke, is taking blood thinners, or is a child or teenager (risk of Reye's syndrome).
\item Conscious: give fast-acting sugar (glucose tablets, juice, or sugary drink) and stay with them; recheck symptoms. Unconscious: call emergency services, do NOT give food or drink (choking risk), place them in the recovery position, and monitor breathing.
\item Suspect opioid overdose (pinpoint pupils, slow or absent breathing, unresponsiveness). Administer naloxone (Narcan) if available and you are trained, one dose into the nose or muscle, repeating every 2--3 minutes if there is no response, and give rescue breathing if needed.
\end{enumerate}

\section*{Chapter 7: Injury Management}

\begin{enumerate}[leftmargin=*]
\item Superficial: red, dry, and painful, affecting the top layer of skin (no blisters). Partial-thickness: blisters, moist, and very painful. Full-thickness: white, leathery, or charred, often painless because nerve endings are destroyed.
\item The Rule of Nines estimates the percentage of body surface area burned, guiding decisions about fluid needs and the need for emergency care. In an adult: head 9\%, each arm 9\%, front trunk 18\%, back trunk 18\%, each leg 18\%, and perineum 1\%.
\item Cool the burn with running water as soon as possible (within 3 hours) for about 20 minutes. Do not use ice or ice water, and avoid prolonged cooling of large burns (risk of hypothermia).
\item Clothing stuck to a burn is adhered to damaged skin; pulling it off tears skin, increases pain, and raises infection risk. Cut around it and cool through/around it, leaving the adhered portion in place.
\item RICE: Rest, Ice, Compression, Elevation. It is appropriate for sprains, strains, and soft-tissue injuries without signs of fracture. Apply for the first 24--48 hours.
\item A fall that leaves someone unable to move their legs suggests possible spinal injury; moving them risks permanent paralysis. Call emergency services, keep them still, support the head and neck in line, and wait for professional rescuers.
\end{enumerate}

\section*{Chapter 8: Environmental First Aid}

\begin{enumerate}[leftmargin=*]
\item Heat exhaustion: move to a cool place, rest, rehydrate, and cool with cloths/fans; it usually responds to these measures. Heat stroke: a life-threatening emergency -- call emergency services and cool the person aggressively (cool or iced water, cool packs to neck/armpits/groin) even before they arrive.
\item Gentle handling avoids triggering dangerous arrhythmias. Warming extremities directly drives cold blood back to the core ("afterdrop"), dropping core temperature further and risking cardiac arrest. Warm the core first and remove wet clothing.
\item Rinse the area with seawater. Hot water (up to 45\textdegree C) or a cold pack may relieve the pain. Do not use fresh water, and use vinegar only for box jellyfish stings.
\item Seek care for flu-like symptoms (fever, chills, headache, muscle aches), inability to remove the whole tick, signs of serious tick-borne disease (facial droop, severe headache, stiff neck, joint swelling, palpitations), or a known attachment of 24 hours or more.
\item Descend immediately. Descent is mandatory at the first signs of severe altitude sickness, high-altitude pulmonary edema, or cerebral edema; it is the only reliable treatment. Use supplemental oxygen if available.
\item Check responsiveness and breathing immediately and start CPR if needed -- victims of lightning are not "electrified." Respiratory arrest is common, so begin rescue breaths/CPR promptly and use an AED if available.
\end{enumerate}

\section*{Chapter 9: Special Populations}

\begin{enumerate}[leftmargin=*]
\item Infants and young children have a higher surface-area-to-body-mass ratio, immature temperature regulation, faster fluid turnover, and cannot communicate thirst or symptoms clearly.
\item Place her on her left side (left lateral position) -- this lifts the uterus off the vena cava, improving blood flow back to the heart and blood pressure for both mother and baby.
\item Weak bones that fracture easily (osteoporosis), medications such as blood thinners that increase bleeding risk, reduced balance and slower reaction time, and slower healing with greater risk of complications.
\item Face the person, get their attention first, speak clearly at a moderate pace, and use gestures or writing if needed. Ensure hearing aids are in place and working.
\item Check the wrist, then the neck, and then the ankle. Medical alert IDs may also be worn on a necklace or bracelet.
\item In late pregnancy the uterus pushes organs upward, so hand placement for compressions is moved slightly higher on the sternum, and the person is tilted onto the left side (or the uterus manually displaced) to improve blood return during CPR.
\end{enumerate}

\section*{Chapter 10: Disaster Response and Preparedness}

\begin{enumerate}[leftmargin=*]
\item START stands for Simple Triage and Rapid Treatment. Its four outcome categories are Immediate (Red -- treatment now), Delayed (Yellow -- can wait), Minor (Green -- walking wounded), and Deceased/Expectant (Black).
\item The quiet, unresponsive person. In START, the loud walking person is marked Green and asked to move to a designated area, while the non-responding person gets rapid assessment and is prioritized as Immediate (Red).
\item In an earthquake, remain in your car (it protects from falling debris); if indoors, drop, cover, and hold on -- get under sturdy furniture, protect your head and neck, and stay away from windows.
\item Call 112 (or the local emergency number), get out immediately, stay low under the smoke, feel doors with the back of your hand before opening, close doors behind you, and go to the family meeting point. Never re-enter a burning building.
\item Local phone lines and cell towers are often overwhelmed or damaged in a disaster, but long-distance lines may still work. One out-of-area contact lets everyone call one person to report they are safe.
\item Water (at least 1 gallon per person per day), non-perishable food, a first aid kit, flashlight with spare batteries, battery/ hand-crank radio, medications and personal supplies, sanitation items, cash, copies of important documents, and items for infants/pets as needed -- enough to last 72 hours.
\end{enumerate}

\section*{Chapter 11: Wilderness and Extended Care First Aid}

\begin{enumerate}[leftmargin=*]
\item Priorities shift from stabilize-and-hand-over to survival: protect the person from the environment, secure shelter, warmth, and water, monitor and reassess repeatedly, and make conservative evacuation decisions. Avoid interventions you cannot sustain or monitor.
\item Rewarm in the field only if you can keep the extremity continuously warm and there is no risk of refreezing. If refreezing is possible, keep the tissue frozen until definitive care -- refreezing causes far more damage than staying frozen. Never rub the area.
\item Tape or butterfly closures trap bacteria inside deep or dirty wounds, greatly increasing the risk of serious infection. Leave the wound open and seek medical care.
\item (1) Is the person's life in danger where they are? (2) Can they be moved without serious harm (suspected spinal injury, open fracture, severe head injury)? (3) How far, how long, and how many people does the move require? (4) Is waiting with shelter, warmth, supplies, and a signal safer than a risky move?
\item "Three" of anything signals distress (three whistle blasts, three shouts, three fires in a triangle, three flashes of light). The buddy system ensures someone knows your route and expected return time and will raise the alarm if you are overdue -- it saves more lives than any equipment.
\end{enumerate}

\section*{Chapter 12: First Aid Practice Scenarios}

\begin{enumerate}[leftmargin=*]
\item Gasping (agonal) breaths are not effective breathing and are common in cardiac arrest. Treating them as breathing delays CPR -- begin compressions immediately.
\item Removing a blood-soaked dressing tears away the forming clot and restarts severe bleeding. Add additional dressings on top instead.
\item Time-critical treatments (e.g., clot-busting drugs for stroke) must be given within a narrow window from the moment symptoms began. The recorded onset time determines eligibility.
\item Nothing placed in the mouth can stop a seizure; it risks choking, broken teeth, and injury. Protect the person, clear hazards, cushion the head, and time the seizure.
\item A fall with inability to move the legs suggests possible spinal injury; moving or lifting could cause permanent paralysis. Keep the person still, maintain manual neck stabilization, and wait for professional rescuers.
\item Naloxone wears off in 30--90 minutes and a second overdose can occur as it wears off. The person must be monitored by medical personnel even if they wake up.
\item Rapid cooling is the top priority because the high body temperature damages organs continuously; every minute of delay increases injury. Cooling must begin immediately, even during transport.
\item The embedded object may be plugging the wound and controlling bleeding. Removing it can trigger severe, uncontrolled bleeding and further tissue damage.
\item Water in the lungs can cause delayed respiratory failure (secondary drowning) hours after the rescue, even when the child seemed fine. Any submersion with new symptoms requires urgent medical evaluation.
\end{enumerate}

\section*{Appendix C: Emergency Report Form and Contact Numbers}

\begin{enumerate}[leftmargin=*]
\item To give responders an accurate timeline. Treatment decisions and response priority depend on time -- the time of collapse, the time CPR started, the time symptoms began, and the time EMS was called.
\item Your exact location, what happened, how many people are involved, and the condition of the person (responsive? breathing?) plus what is being done. Do not hang up until the dispatcher says to.
\item A medical alert ID carries critical medical information (conditions, allergies, emergency contact) that responders rely on. Removing it loses that information -- tell responders exactly what it says instead.
\end{enumerate}
